\lecture{Древнерусское искусство}{lecture-01}

{
  \title{История Русского Искусства}
  \subtitle{Лекция 1: Древнерусское искусство}
  \date{4 сентября 2026}

  \begin{frame}[plain,noframenumbering]
    \titlepage
  \end{frame}
}

\sectionimage{sirin.png}
\section{Искусство восточных славян}

\begin{frame}{Расселение древних славян}
  \tallfigure[]
    {lectures/01-ancient/images/maps/map1.jpg}
    {Три ветви расселения древних славян к VI-IX вв.}

  Расселение происходило по 3 основным направлениям:
  \begin{itemize}
    \item на юг, на Балканский полуостров; 
    \item на запад, близ междуречья Одера и Эльбы; 
    \item на восток и север, по Восточноевропейской равнине.
  \end{itemize}
\end{frame}

\begin{frame}{Искусство восточных славян}
  Дописьменное искусство ранних славян это прежде всего металлопластика
  и языческая пластика. Основные источники:
  \begin{itemize}
    \item подвески-обереги с геометрическим и солярным орнаментом;
    \item прикладное кузнечное искусство: фибулы, браслеты, височные кольца;
    \item украшения вятичей: височные кольца, гривна, браслеты, кресты;
    \item сохранившиеся сооружения: могильники, святилища, жилищная и крепостная архитектура;
    \item кумиры и идолы, т.е. каменные и деревянные изваяния божеств.
  \end{itemize}
\end{frame}

\subsection{Мартыновский клад}

\begin{frame}{Мартыновский клад}
  \begin{itemize}
    \item VI век, село Мартыновка в устье реки Рось.
    \item Один из ключевых комплексов раннеславянской
    металлопластики.
    \item Находка восьми литых фигурок людей и коней.
    \item Серебряные нашивные бляшки, отлитые в формах.
    \item Обработка деталей чеканкой.
    \item Уркашение позолотой.
    \item Служили оберегами и амулетами.
    \item Реалистическая трактовка деталей и сильная стилизация, чистая орнаментальность напоминают скифо-сарматский звериный стиль.
  \end{itemize}
\end{frame}

\begin{frame}{Мартыновский клад}
  \widefigure[]
    {lectures/01-ancient/images/maps/map2.png}
    {}
\end{frame}

\begin{frame}{Мартыновский клад}
  \tallfigure[фото: Völkerwanderer; Британский музей;
    \href{https://commons.wikimedia.org/wiki/File:Slavic_silver_hoard.JPG}{Wikimedia Commons}. CC0]
    {lectures/01-ancient/images/martynivka-hoard/figure-03-hoard-british-museum.jpg}
    {Часть Мартыновского клада, ок.\ 550--650~гг. Серебряные бляшки,
     зооморфные накладки и украшения. Британский музей, Лондон}
\end{frame}

\begin{frame}{«Танцующий человек»}
  \tallfigure[\href{https://commons.wikimedia.org/wiki/File:Anty_bljashka.jpg}{Wikimedia Commons}.
    Общественное достояние]
    {lectures/01-ancient/images/martynivka-hoard/figure-01-anthropomorphic-plaque.jpg}
    {Серебряная нашивная бляшка из Мартыновского клада.
     VI--VII~вв. Антропоморфная фигурка в «танцующей» позе;
     отверстия у локтей --- для крепления к одежде или сбруе}
\end{frame}

\begin{frame}{Пара бляшек}
  \widefigure[Музей исторических драгоценностей Украины, Киев;
    \href{https://commons.wikimedia.org/wiki/File:Anty,_bljashky.jpg}{Wikimedia Commons}.
    Общественное достояние]
    {lectures/01-ancient/images/martynivka-hoard/figure-02-dancing-men-pair.jpg}
    {Пара антропоморфных серебряных бляшек («танцующие люди»).
     Высота 7,6 и 7,5~см. Мартыновский клад, находка 1909~г.}
\end{frame}

\subsection{Амулеты и обереги}

\begin{frame}{Амулеты и обереги}
  \begin{itemize}
    \item Бронзовые и серебряные подвески из кладов и погребений.
    \item Геометрический и солярный орнамент; лунницы; зооморфные «коники».
    \item Обереги --- тесная связь украшения с языческим мироощущением.
  \end{itemize}
\end{frame}

\begin{frame}{Солярный амулет}
  \tallfigure[фото: Чаховіч Уладзіслаў; Могилёвский музей;
    \href{https://commons.wikimedia.org/wiki/File:\%D0\%9F\%D0\%B0\%D0\%B4\%D0\%B2\%D0\%B5\%D1\%81\%D0\%BA\%D0\%B0\%E2\%80\%93\%D0\%B0\%D0\%BC\%D1\%83\%D0\%BB\%D0\%B5\%D1\%82_\%D0\%B7_\%D1\%81\%D0\%B0\%D0\%BB\%D1\%8F\%D1\%80\%D0\%BD\%D1\%8B\%D0\%BC_\%D1\%81\%D1\%96\%D0\%BC\%D0\%B2\%D0\%B0\%D0\%BB\%D0\%B0\%D0\%BC.jpg}{Wikimedia Commons}.
    Общественное достояние]
    {lectures/01-ancient/images/slavic-amulets/solar-amulet-pendant.jpg}
    {Подвеска-амулет с солярным символом.
     Бронза, литьё. Конец X --- начало XI~в.}
\end{frame}

\begin{frame}{Лунница}
  \tallfigure[фото: Чаховіч Уладзіслаў; Могилёвский музей;
    \href{https://commons.wikimedia.org/wiki/File:\%D0\%9F\%D0\%B0\%D0\%B4\%D0\%B2\%D0\%B5\%D1\%81\%D0\%BA\%D0\%B0-\%D0\%BB\%D1\%83\%D0\%BD\%D0\%BD\%D1\%96\%D1\%86\%D0\%B0_\%D0\%9A\%D0\%B0\%D0\%BD\%D0\%B5\%D1\%86_XI_-_\%D0\%BF\%D0\%B0\%D1\%87\%D0\%B0\%D1\%82\%D0\%B0\%D0\%BA_XII_\%D1\%81\%D1\%82..jpg}{Wikimedia Commons}.
    Общественное достояние]
    {lectures/01-ancient/images/slavic-amulets/lunnitsa-pendant.jpg}
    {Лунница --- подвеска в форме полумесяца.
     Серебро, литьё. Конец XI --- начало XII~в. Длина 2,7~см}
\end{frame}

\begin{frame}{Амулет «коник»}
  \tallfigure[\href{https://commons.wikimedia.org/wiki/File:\%D0\%90\%D0\%BC\%D1\%83\%D0\%BB\%D0\%B5\%D1\%82_\%C2\%AB\%D0\%9A\%D0\%BE\%D0\%BD\%D1\%96\%D0\%BA\%C2\%BB_1.jpg}{Wikimedia Commons}.
    Общественное достояние]
    {lectures/01-ancient/images/slavic-amulets/horse-amulet-konik.jpg}
    {Бронзовый зооморфный амулет-подвеска («коник»)
     с круговым солярным орнаментом.
     Типичный оберег женского убора}
\end{frame}

\begin{frame}{Языческие амулеты}
  \tallfigure[фото: Лапоть; Владимирский исторический музей;
    \href{https://commons.wikimedia.org/wiki/File:Amulets_-_Vladimir_Historical_Museum.jpg}{Wikimedia Commons}. CC0]
    {lectures/01-ancient/images/slavic-amulets/amulets-vladimir-museum.jpg}
    {Языческие амулеты XI --- начала XIII~в.:
     клыки и когти, птицевидные подвески, бусы.
     Владимирский исторический музей}
\end{frame}

\subsection{Прикладное искусство}

\begin{frame}{Прикладное искусство}
  Прикладное искусство распрастраненный вид деятельности. Изучен лучше всего и являлся наиболее стойким оружием в борьбе против враждебной христианской идеологии. 
  
  Кузнечное и ювелирное ремесло восточных славян являлось основным видом деятельности и находилось в центре культурной жизни племён:
  \begin{itemize}
    \item фибулы, подвески, браслеты, височные кольца;
    \item предметы быта, игрушки, посуда;
    \item орнамент и форма украшений связаны с языческим мироощущением
          и племенной принадлежностью.
  \end{itemize}
\end{frame}

\begin{frame}{Фибулы}
  \tallfigure[фото: Daderot; Morgan Library \& Museum, Нью-Йорк;
    \href{https://commons.wikimedia.org/wiki/File:Pair_of_radiate-head_bow_brooches,_Slavic,_1_of_2,_c._600-650_AD,_copper_alloy,_gilding_-_Morgan_Library_\%26_Museum_-_New_York_City_-_DSC06618.jpg}{Wikimedia Commons}.
    Общественное достояние]
    {lectures/01-ancient/images/slavic-applied-art/slavic-fibulae-morgan-library.jpg}
    {Пальчатая фибула. Славянская работа,
     ок.\ 600--650~гг. Медный сплав с позолотой}
\end{frame}

\begin{frame}{Височные кольца}
  \widefigure[фото: Лапоть; Музей Москвы;
    \href{https://commons.wikimedia.org/wiki/File:Temple_rings_-_11-12c_-_Museum_of_Moscow.jpg}{Wikimedia Commons}. CC0]
    {lectures/01-ancient/images/slavic-applied-art/temple-rings-museum-of-moscow.jpg}
    {Височные кольца XI--XII~вв.: щитковые, бусинные,
     пяти- и семилопастные, семилучевые. Музей Москвы}
\end{frame}

\begin{frame}{Народные украшения}
  Народные мастера все чаще создают подлинные произведения искусства, чьи декоративные элементы тесно связаны с языческим мировозрением и символикой.
  \begin{itemize}
    \item Ажурное литье
    \item Гармоничные формы, геометрический орнамент
    \item Цветная выемчатая эмаль свидетельствует об уже высоком иссустве высокого уровня.

  \end{itemize}
\end{frame}

\begin{frame}{Украшения вятичей}
  \widefigure[фото: Лапоть; Тульский областной краеведческий музей;
    \href{https://commons.wikimedia.org/wiki/File:Jewellery_-_Vyatichi_-_Tula_regional_museum_of_local_lore.jpg}{Wikimedia Commons}. CC0]
    {lectures/01-ancient/images/slavic-applied-art/temple-rings-vyatichi-tula-museum.jpg}
    {Женский убор вятичей: височные кольца, гривна, браслеты, кресты}
\end{frame}

\begin{frame}{Браслеты}
  \widefigure[фото: Лапоть; Музей археологии Москвы;
    \href{https://commons.wikimedia.org/wiki/File:Bracelets_-_Rus\%27_12-13c_-_Museum_of_Archaeology_of_Moscow.jpg}{Wikimedia Commons}. CC0]
    {lectures/01-ancient/images/slavic-applied-art/bracelets-rus-museum-archaeology-moscow.jpg}
    {Бронзовые браслеты. Русь, XII--XIII~вв.}
\end{frame}

\begin{frame}{Украшения дреговичей}
  \widefigure[фото: Лапоть; Государственный исторический музей, Москва;
    \href{https://commons.wikimedia.org/wiki/File:Jewellery_Dregoviches_GIM.jpg}{Wikimedia Commons}. CC0]
    {lectures/01-ancient/images/slavic-applied-art/jewellery-dregoviches-gim.jpg}
    {Вещи дреговичей XI--XIII~вв.: височные кольца, бусы,
     детали костюма. ГИМ}
\end{frame}

\subsection{Кумиры}

\begin{frame}{Кумиры}
  Подобно палеолитическим венерам, кумиры являлись символами плодородия и защиты от бедствий. 
  \begin{itemize}
    \item Скульптурные изображения из бронзы, дерева, глины, камня с незначительной проработкой деталей.
    \item Детали примитивные, графические или барельефные.
  \end{itemize}
\end{frame}

\begin{frame}{Кумиры}
  \begin{itemize}
    \item Акулининское изваяние (близ Подольска, ныне Домодедовский район)
          --- небольшое известняковое поясное изображение,
          возможно женское божество. Хранится в ГИМ;
          свободной фотографии нет.
    \item Наиболее яркий пример изображения женского божества. Погрудная фигура из известняка. Голова в техинке кругой, обьемной скульптуры. Черты лица просто процарапаны резцом, в профиль не просматриваются.
  \end{itemize}
\end{frame}

\begin{frame}{Псковский идол}
  \tallfigure[фото 1928--1929~гг.;
    \href{https://commons.wikimedia.org/wiki/File:Pskov_idol_1.jpg}{Wikimedia Commons}.
    Общественное достояние]
    {lectures/01-ancient/images/slavic-idols/pskov-idol.jpg}
    {Каменный идол («баба»), найден в 1897~г.\ близ Пскова.
     VIII--X~вв.}
\end{frame}

\begin{frame}{Идолы}
  Культовые изваяния не унифицированы, нет четкой иконографии. Каждый идол имеет свою индивидуальность и символику.
  \begin{itemize}
    \item Были погудными в рост, чаще изображалась одна голова на длинном деревянном или каменном древке-столпе. (Свидетельство средневековых арабских источников)
  \end{itemize}
\end{frame}

\begin{frame}{Сиверский идол}
  \tallfigure[фото: Shakko; Музей крестьянского быта, Ферапонтов монастырь;
    \href{https://commons.wikimedia.org/wiki/File:Siverskiy_idol.jpg}{Wikimedia Commons}.
    Общественное достояние]
    {lectures/01-ancient/images/slavic-idols/siverskiy-idol.jpg}
    {Сиверский идол. IV--IX~вв., высота ок.\ 1~м.
     Деревня Сиверово}
\end{frame}

\subsection{Збручский идол}

\begin{frame}{Збручский идол}
  \begin{itemize}
    \item IX--X~вв. Четырёхгранный известняковый столб, высота 257~см.
    \item Найден на холме над рекой Збруч,
          на границе земель волынян, белых хорватов, бужан и тиверцев.
    \item Три яруса: небо --- земля --- подземный мир.
    \item Археологический музей в Кракове.
  \end{itemize}
\end{frame}

\begin{frame}{Збручский идол}
  \tallfigure[Археологический музей в Кракове;
    оцифровка: Małopolska's Virtual Museums;
    \href{https://commons.wikimedia.org/wiki/File:Svetovid_\%E2\%80\%93_Zbruch_Idol_(MAK-63).jpg}{Wikimedia Commons}.
    Общественное достояние]
    {lectures/01-ancient/images/zbruch-idol/zbruch-idol-museum-photo.jpg}
    {Збручский идол (Святовид). Известняк, IX--X~вв.}
\end{frame}

\begin{frame}{В экспозиции}
  \tallfigure[фото: Silar, 2013; Археологический музей в Кракове;
    \href{https://commons.wikimedia.org/wiki/File:O98_Idol_von_Sbrutsch_mit_Darstellung_von_Unterwelt,_Erde_und_des_Himmels,_zirka_10._Jh._n._Chr..JPG}{Wikimedia Commons}.
    CC~BY-SA~3.0]
    {lectures/01-ancient/images/zbruch-idol/zbruch-idol-krakow-museum-45.jpg}
    {Три яруса космологии: небо, земля, подземный мир}
\end{frame}

\begin{frame}{Четыре грани}
  \widefigure[рисунок: Прозоров Сергей;
    \href{https://commons.wikimedia.org/wiki/File:\%D0\%A1\%D1\%82\%D0\%BE\%D1\%80\%D0\%BE\%D0\%BD\%D1\%8B_\%D0\%B7\%D0\%B1\%D1\%80\%D1\%83\%D1\%87\%D1\%81\%D0\%BA\%D0\%BE\%D0\%B3\%D0\%BE_\%D0\%B8\%D0\%B4\%D0\%BE\%D0\%BB\%D0\%B0.jpg}{Wikimedia Commons}.
    Общественное достояние]
    {lectures/01-ancient/images/zbruch-idol/zbruch-idol-four-sides-diagram.jpg}
    {Четыре грани Збручского идола: божества верхнего яруса
     (рог, кольцо, конь и меч) и поддерживающая фигура нижнего мира}
\end{frame}

\subsection{Капища}

\begin{frame}{Славянская языческая религия}
  \begin{itemize}
    \item Бог Род (Сварог, Святовит/д) --- плодородие, природа, жизнь, повелитель молний и дождя.
    \item Солнечние боги Даждьбог, Хорс, Ярила, Перун, Велес.
    \item Нижняя ступень --- Ветры, Русалки, Рожаницы. Вместе с почитаемыми богами предками даровали изобилие.
  \end{itemize}

  Позднее на первом месте оказывается Перун --- бог феодальной верхушки общества, князей и дружинников.
\end{frame}

\begin{frame}{Славянская языческая религия}
  <<И нача княжити Володимер в Киеве един и постави кумиры на холму вне двора теремного: Перуна древяна, а главу его сребрену, а ус злат, и хЬрса, Дажьбога и Стрибога, и Симаргла, и Мокощь>>

  Летопись конца X века на княжение Владимира Святославича.
\end{frame}

\begin{frame}{Капище}
  \begin{itemize}
    \item Языческое святилище, культовое сооружение
          разнообразных архитектурных форм.
    \item Открытые площадки с рвом и идолом; деревянные храмы.
    \item Капище у реки Гнилопять под Житомиром (Шумск):
          раскопки И.~П.~Русановой, IX~в.;
          свободной фотографии раскопа нет.
  \end{itemize}
\end{frame}

\begin{frame}{Святилище Перуна на Перыни}
  \tallfigure[рисунок: Sterndmitri, по В.~В.~Седову;
    \href{https://commons.wikimedia.org/wiki/File:PerunSchrineScheme.png}{Wikimedia Commons}.
    CC~BY-SA~4.0]
    {lectures/01-ancient/images/kapishche/peryn-shrine-scheme.png}
    {План капища Перуна на Перыни близ Новгорода, X~в.:
     круглая площадка (ок.\ 21~м), восьмилепестковый ров
     с кострищами, яма от идола в центре}
\end{frame}

\begin{frame}{Реконструкция Перыни}
  \widefigure[рисунок: Віщун, 2021; по плану Седова;
    \href{https://commons.wikimedia.org/wiki/File:Shrine_of_Perun_in_Peryn.jpg}{Wikimedia Commons}.
    CC~BY-SA~4.0]
    {lectures/01-ancient/images/kapishche/peryn-shrine-reconstruction.jpg}
    {Реконструкция открытого святилища: деревянный идол на холме,
     круговой ров, ритуальные костры}
\end{frame}

\begin{frame}{Культовая архитектура}
  \begin{itemize}
    \item Конструкция и типология помещений для поклонения идолам и жертвоприношений весьма разнообразны и плохо изучены.
  \end{itemize}
\end{frame}

\begin{frame}{Деревянный храм}
  \widefigure[фото: Wolfgang Sauber; музей Groß Raden;
    \href{https://commons.wikimedia.org/wiki/File:Freilichtmuseum_Gro\%C3\%9F_Raden_-_Tempel_3.jpg}{Wikimedia Commons}.
    CC~BY-SA~4.0]
    {lectures/01-ancient/images/kapishche/gross-raden-slavic-temple.jpg}
    {Реконструкция славянского деревянного храма X~в.
     Радоним / Groß Raden, Мекленбург}
\end{frame}

\subsection{Западнославянские храмы}

\begin{frame}{Аркона и Ретра}
  \begin{itemize}
    \item Храм Святовита в Арконе (мыс Аркона, Рюген):
          четырёхликий кумир с турьим рогом.
    \item Храм Радогоста (Радегаста) в Ретре ---
          святилище лютичей; точное место не установлено.
  \end{itemize}
\end{frame}

\begin{frame}{Аркона}
  \widefigure[рисунок: Віщун, 2023;
    \href{https://commons.wikimedia.org/wiki/File:Arkona_artistic_reconstruction.jpg}{Wikimedia Commons}.
    CC~BY-SA~4.0]
    {lectures/01-ancient/images/west-slavic-temples/arkona-artistic-reconstruction.jpg}
    {Реконструкция Яромарсбурга на мысе Аркона:
     валы, жилая застройка и храм Святовита на обрыве}
\end{frame}

\begin{frame}{Святовит}
  \tallfigure[рисунок: Bornholm, 2023; по Саксону Грамматику;
    \href{https://commons.wikimedia.org/wiki/File:Svetovid_\%C5\%9Awi\%C4\%99towid_wg_Saxo_Grammaticus.jpg}{Wikimedia Commons}.
    CC~BY-SA~4.0]
    {lectures/01-ancient/images/west-slavic-temples/svantevit-according-to-saxo.jpg}
    {Кумир Святовита по описанию Саксона Грамматика:
     четыре лица, в руке --- турий рог}
\end{frame}

\begin{frame}{Ретра}
  \widefigure[рисунок: Ronny Krüger, 2011;
    \href{https://commons.wikimedia.org/wiki/File:RekonstruktionRethra.jpg}{Wikimedia Commons}. CC0]
    {lectures/01-ancient/images/west-slavic-temples/rethra-reconstruction.jpg}
    {Гипотетическая реконструкция Ретры, X~в.:
     островное городище, центральный храм Радогоста}
\end{frame}

\subsection{Жилищная и крепостная архитектура}

\begin{frame}{Светская архитектура}
  \begin{itemize}
    \item Основной источник это археологические раскопки и сведения. Архитектура мало изучена.
    \item Планировка деревоземляных, глинобитных и каменных сооружений, южных полуземлянок и северных срубов.
  \end{itemize}
\end{frame}

\begin{frame}{Полуземлянка}
  \widefigure[фото: Silar, 2016; музей Trzcinica;
    \href{https://commons.wikimedia.org/wiki/File:02016_0036_Rekonstruktion_eines_slawischen_Grubenhauses_im_Freilichtmuseum_Trzcinica,_Beskiden.jpg}{Wikimedia Commons}.
    CC~BY-SA~4.0]
    {lectures/01-ancient/images/slavic-architecture/figure-01-semidugout-trzcinica.jpg}
    {Реконструкция славянской полуземлянки, ок.\ 800--900~гг.}
\end{frame}

\begin{frame}{Интерьер полуземлянки}
  \tallfigure[фото: Silar;
    \href{https://commons.wikimedia.org/wiki/File:020210828_Typical_equipment_of_a_Slavic_house_from_the_period_of_their_migrations,_typical_half-dugout,_Slavic_settlement,_Archaeology.jpg}{Wikimedia Commons}.
    CC~BY-SA~4.0]
    {lectures/01-ancient/images/slavic-architecture/figure-02-semidugout-interior.jpg}
    {Типичное убранство полуземлянки эпохи расселения славян}
\end{frame}

\begin{frame}{Сруб}
  \widefigure[фото: Wolfgang Sauber, 2014; музей Groß Raden;
    \href{https://commons.wikimedia.org/wiki/File:Freilichtmuseum_Gro\%C3\%9F_Raden_-_Blockhaus.jpg}{Wikimedia Commons}.
    CC~BY-SA~4.0]
    {lectures/01-ancient/images/slavic-architecture/figure-03-gross-raden-log-cabin.jpg}
    {Реконструкция западнославянского сруба с соломенной кровлей, X~в.}
\end{frame}

\begin{frame}{Городище}
  \widefigure[фото: Wolfgang Sauber, 2014; музей Groß Raden;
    \href{https://commons.wikimedia.org/wiki/File:Freilichtmuseum_Gro\%C3\%9F_Raden_-_Wall_und_Graben.jpg}{Wikimedia Commons}.
    CC~BY-SA~4.0]
    {lectures/01-ancient/images/slavic-architecture/figure-04-gross-raden-palisade.jpg}
    {Вал, ров и частокол IX--X~вв. --- крепостная архитектура городища}
\end{frame}

\subsection{Чёрная могила}

\begin{frame}{Турий рог}
  \begin{itemize}
    \item Большой турий рог из княжеского кургана «Чёрная могила»
          в Чернигове. IX--X~вв.
    \item Серебряная оковка: сцена охоты / боя.
    \item В литературе связывают с былиной об Иване Годиновиче.
    \item Государственный исторический музей, Москва.
    \item Изображены вещая птица орёл, слева длиннокосая девушка с луком и колчаном, справа бог войны, в руке меч.
  \end{itemize}
\end{frame}

\begin{frame}{Курган}
  \widefigure[фото: И.~В.~Тестов;
    \href{https://commons.wikimedia.org/wiki/File:\%D0\%9A\%D1\%83\%D1\%80\%D0\%B3\%D0\%B0\%D0\%BD_\%D0\%A7\%D1\%91\%D1\%80\%D0\%BD\%D0\%B0\%D1\%8F_\%D0\%BC\%D0\%BE\%D0\%B3\%D0\%B8\%D0\%BB\%D0\%B0.jpg}{Wikimedia Commons}.
    CC~BY-SA~3.0]
    {lectures/01-ancient/images/chernaya-mogila/kurgan-chernaya-mogila.jpg}
    {Курган «Чёрная могила» в Чернигове.
     Раскопки Д.~Я.~Самоквасова, 1872--1873}
\end{frame}

\begin{frame}{Большой турий рог}
  \widefigure[фото: Лапоть; Государственный исторический музей, Москва;
    \href{https://commons.wikimedia.org/wiki/File:Rhyton_Black_Grave_GIM_Zlatnik.jpg}{Wikimedia Commons}. CC0]
    {lectures/01-ancient/images/chernaya-mogila/large-rhyton-gim.jpg}
    {Турий рог с серебряной золочёной оковкой. X~в. ГИМ}
\end{frame}

\begin{frame}{Оковка}
  \widefigure[\href{https://commons.wikimedia.org/wiki/File:Rhyton_Black_Grave_01.jpg}{Wikimedia Commons}.
    Общественное достояние]
    {lectures/01-ancient/images/chernaya-mogila/overlay-chase-detail.jpg}
    {Деталь чеканки большого рога: лучники и птица.
     Сцену соотносят с былиной об Иване Годиновиче}
\end{frame}
